\chapter*{ABSTRAK}
\addcontentsline{toc}{chapter}{ABSTRAK}

\begin{center}

    {\Large\bfseries Pengembangan Aplikasi Pembelajaran Algoritma dan Struktur Data Berbasis \textit{Large Language Model}}\\
    oleh \\
    \vspace{0.4cm}
    Farah Aulia\\
    18222096\\
\end{center}


\begin{spacing}{1.0}
Penelitian ini mengembangkan Teman Belajar ASD, aplikasi pembelajaran Algoritma dan Struktur Data (ASD) berbasis \textit{web} yang memanfaatkan \textit{Large Language Model} (LLM) sebagai teman belajar digital. Aplikasi menerapkan pendekatan \textit{guided learning} dan \textit{prompt engineering} untuk menghasilkan materi, contoh implementasi, soal latihan adaptif, dan ringkasan yang relevan serta konsisten dengan kurikulum ASD.

ASD merupakan mata kuliah fundamental yang kerap sulit dipahami mahasiswa karena materinya abstrak dan multidimensional, sementara metode pembelajaran konvensional belum menyediakan umpan balik yang cepat dan adaptif di luar jam kuliah. Kondisi ini mendorong kebutuhan akan media belajar mandiri yang lebih interaktif dan responsif.

Aplikasi dikembangkan menggunakan kerangka \textit{Software Development Life Cycle} (SDLC) meliputi tahap perencanaan, analisis, desain, implementasi, dan pengujian, dengan \textit{frontend} Next.js dan React, \textit{backend} FastAPI, basis data PostgreSQL, serta integrasi LLM melalui OpenAI API model GPT-4o-mini. Selain alur \textit{guided learning} utama, aplikasi menyediakan fitur evaluasi soal mandiri dan asisten percakapan kontekstual. Evaluasi dilakukan melalui pengujian \textit{black box}, \textit{usability testing} menggunakan \textit{System Usability Scale} (SUS), serta evaluasi kualitas konten LLM menggunakan rubrik ahli.

Hasil pengujian \textit{black box} menunjukkan seluruh 18 skenario uji dari 14 fitur fungsional dinyatakan lulus. \textit{Usability testing} dengan 5 responden menghasilkan skor SUS rata-rata 93,0 (kategori \textit{Excellent}). Evaluasi kualitas konten LLM terhadap 15 soal referensi pada tiga topik ASD memperoleh rata-rata kesesuaian 97,33\%, melampaui indikator keberhasilan minimal 70\%. Hasil ini membuktikan Teman Belajar ASD mampu menghadirkan pengalaman belajar ASD yang terstruktur, interaktif, dan konsisten dengan kurikulum.

\vspace{0.5em}
\noindent\textbf{Kata kunci:} Algoritma dan Struktur Data, \textit{Large Language Model}, \textit{Guided Learning}, \textit{Prompt Engineering}, Aplikasi Pembelajaran Berbasis \textit{Web}.
\end{spacing}
